\section{Discussion}
\label{sect:discussion}

\setlength{\fboxsep}{1.2pt}
\textbf{Future directions for sustainable AI agent serving.}
While the primary objective of our work is to raise awareness within the community about the broader system-level implications of deploying agents—particularly the immense infrastructural costs associated with them—we also highlight several promising directions that we believe will be critical for sustainably serving agentic systems.
First, conventional \textit{model-level optimizations} such as quantization~\cite{lin2024awq,xiao2023smoothquant,frantar2022gptq}, distillation~\cite{hinton2015distilling,muralidharan2024compact,sreenivas2024llm,guo2025deepseek}, sparse architectures~\cite{fedus2022switch,yuan2025native,liang2024mixture}, and adaptive model routing~\cite{panda2025adaptive,ding2025best} will remain essential for reducing computational and memory demands. For example, constructing multi-agent systems that combine a heterogeneous mix of small and large language models (SLMs and LLMs), and dynamically selecting the appropriate model depending on the agent’s role and task significance (e.g., planning vs. acting), can substantially reduce both operational cost and latency—an approach also advocated by~\cite{belcak2025small}.
Second, for AI agents that are not strictly latency-sensitive, \textit{carbon-aware computing}~\cite{radovanovic2022carbon,li2023clover,gsteiger2024caribou} that migrates parts of the execution to compute instances incurring lower carbon intensity or electricity costs can provide both environmental and economic benefits.
Finally, \textit{adaptive scaling strategies}~\cite{snell2024scaling,qu2025optimizing,yang2025towards} that dynamically adjust compute resources based on task difficulty and importance enable agents to allocate GPU resources more efficiently, avoiding over-provisioning while maintaining quality of service.
Together, these directions highlight exciting opportunities for improving the efficiency, sustainability, and scalability of serving AI agent workloads.


\textbf{Agent serving under SLA constraints.}
The industry has only recently begun exploring agentic systems, so there are currently no well-established or widely accepted SLA standards. Consequently, we did not (and realistically could not) conduct our analysis with respect to a specific SLA target. The goal of this work is to characterize energy–efficiency trade-offs across diverse agent configurations, rather than to optimize for a fixed latency constraint. Considering the current direction of agentic system development, agents are generally allowed to spend additional computation time to achieve higher reasoning quality. Our analysis intentionally examines this behavior to highlight the inefficiency of unrestricted test-time scaling and to motivate more energy-aware SLA design in future agent deployments. A detailed exploration of efficient agent serving under SLA constraints is left as future work.

